This chapter uses a compact, chapter-local notation policy: symbols are reintroduced at first use in each chapter when needed, so that definitions remain close to their point of application. When a symbol is reused with a specialized meaning in a later chapter, that meaning is stated explicitly in the local context.

Throughout the book, we write
\[
\mathbb{B} \coloneqq \{0,1\}
\]
for the Boolean domain, and
\[
\mathbb{B}^n
\]
for the set of bitstrings of length $n$. The notation
\[
\mathrm{Bool}(n)
\]
denotes the set of all $n$-ary Boolean functions, that is, maps $f : \mathbb{B}^n \to \mathbb{B}$.

We use the following conventions for algebraic and compositional structure.
\begin{itemize}
  \item $\mathcal{O}$ denotes an operad, and $\mathcal{O}(n)$ denotes the collection of operations of arity $n$.
  \item $\Sigma_n$ denotes the symmetric group on $n$ letters.
  \item $\circ$ denotes substitution or serial composition.
  \item $\otimes$ denotes parallel composition or monoidal product.
  \item $\mathrm{Tr}$ denotes feedback or trace.
  \item $\mathrm{cl}(\cdot)$ denotes a closure operator.
  \item A \emph{clone} is a set of operations closed under projections and composition.
  \item $G \curvearrowright X$ denotes a group action of $G$ on $X$.
  \item $\mathrm{Aut}(\cdot)$ denotes an automorphism group.
\end{itemize}

Diagrammatic and typing conventions are used consistently across the book. Ports are typed, and when a family of operations carries multiple sorts or colors, we treat it as a colored operadic structure. In such settings, the type of each port is part of the data of the diagram, and well-typed composition is assumed unless stated otherwise.

Error and resolution are measured by a chapter-specific quantity denoted $\varepsilon(r)$, where $r$ indicates the relevant resolution or comparison regime. Unless a chapter overrides it, the default interpretation is normalized Hamming distance. Thus, for $f,g \in \mathrm{Bool}(n)$, we write
\[
\varepsilon_{\mathrm{tt}}(r=n; f,g)
\coloneqq \frac{1}{2^n} \sum_{x \in \mathbb{B}^n} \mathbf{1}[f(x) \neq g(x)].
\]
This quantity is the fraction of inputs on which $f$ and $g$ disagree. When a later chapter requires a different notion of discrepancy---for example, bisimulation distance, metastability-aware error, or another domain-specific metric---that chapter states the override explicitly and uses the same $\varepsilon(r)$ notation with the appropriate interpretation.

Unless otherwise indicated, equality is literal equality of the underlying mathematical objects, composition is read left-to-right in the order specified by the surrounding context, and all diagrams are assumed to be interpreted up to the equivalences declared in the relevant chapter.

For readability, we also adopt the following global conventions. First, subscripts and superscripts are used in the standard mathematical sense unless explicitly redefined. Second, when a chapter introduces a local abbreviation or overloaded symbol, the chapter should restate the intended meaning at the point of use. Third, any statement about equivalence, isomorphism, or invariance is to be read relative to the ambient structure currently under discussion (for example, set-theoretic, algebraic, categorical, or diagrammatic).

These conventions are intended to keep the exposition compact without sacrificing precision: local definitions control local meaning, while the book-wide defaults provide a stable baseline for Boolean, algebraic, and diagrammatic reasoning.
